% QMB_n05_bell_ffgft.tex — Chapter 5: FFGFT Answer to Bell
% Sources: Doc. 230, 190, 353, 364
\chapter{The FFGFT Answer to Bell}
\label{ch:bell_ffgft}

\section{Topological Connection Instead of Action at a Distance}

The current FFGFT position (Doc.~230, May 2026) does not challenge any of the three
Bell assumptions. Instead, it changes the ontological framework:
Entangled particles have no separate carrier, but share
the same torus $T^4/\mathbb{Z}_3$. The correlation is a
topological property of this common carrier.

This means concretely:
\begin{itemize}
  \item \textbf{Locality} holds in the sense of the carrier: No signal
        transmits information faster than light.
  \item \textbf{Realism} holds in the sense that the mode
        configuration exists before the measurement.
  \item \textbf{Freedom of choice} is not restricted.
\end{itemize}
What is given up is the idea that each particle
would have an \emph{independent} carrier. The topological connection
is not action at a distance — it is a geometric fact.\status{B}

\section{Statistical Predictions}

The FFGFT reproduces the Born rule and thus all statistical
predictions of QM completely:\status{B}
\begin{equation}
  P_{\mathrm{FFGFT}}(z=+1) = \tfrac{1+z}{2} = |\alpha|^2 = P_{\mathrm{QM}}(0).
\end{equation}
In particular, the CHSH bound $2\sqrt2$ is reached — not
undershot, as the early damping approaches had predicted. The Bell test dispute over local realism
resolves geometrically: There is no nonlocality, because there is no
separate carrier.

\section{The $\xi$ Deviation}

A small deviation from the Tsirelson bound is derivable from the
$\theta_4$ projection geometry of the measurement act. The
candidate for the elementary deviation per measurement is:\status{S}
\begin{equation}
  \Delta\mathrm{CHSH}_{\mathrm{elem}} = \frac{\xi}{2\pi}
  \approx 2{,}1\times10^{-5}.
\end{equation}
This is the phase cost of a $\theta_4$ cycle in geometric
units (Doc.~190, Precision~12). Cumulatively over $N$ measurement runs,
this deviation grows as:
\begin{equation}
  \Delta\mathrm{CHSH}(N) \approx \xi\,\frac{\ln N}{D_f}\cdot2\sqrt2.
\end{equation}
For $N=73$ this yields $\approx5\times10^{-4}$ — a different
observable than the elementary one, both below NISQ noise
($\sim10^{-2}$).

\textbf{Experimental reach:} Hardware measurements (IBM Heron~r2,
May 2026, Doc.~147) show Bell violation $S=2{,}74$ ($96{,}9\,\%$
of the Tsirelson bound). The $\xi$ effect itself requires
sub-ppm-stable error-corrected qubits — not resolvable with today's
NISQ hardware.\status{K}

\begin{laierspur}
The FFGFT says: The strong quantum correlations are real and
correct — QM is right. What it sees differently: The correlation
does not come from a connection \emph{between} the particles,
but because the particles were never really separate. They sit
on the same geometric carrier like two points on the same
drum.
\end{laierspur}
\quelle{Doc.~190 (Prec.~11, 12), 230, 353, 364, 365}